\begin{thebibliography}{99}
\footnotesize\itemsep=0pt

\bibitem{tlrv}
Alonso~Tarr\'{\i}o L., Jerem\'{\i}as~L\'opez A., P\'erez~Rodr\'{\i}guez M.,
 Vale~Gonsalves M.J., A functorial formalism for quasi-coherent sheaves on a
 geometric stack, \href{https://doi.org/10.1016/j.exmath.2014.12.007}{\textit{Expo. Math.}} \textbf{33} (2015), 452--501,
 \href{https://arxiv.org/abs/1304.2520}{arXiv:1304.2520}.


\bibitem{bdfik}
Ballard M., Deliu D., Favero D., Isik M.U., Katzarkov L., Resolutions in
 factorization categories, \href{https://doi.org/10.1016/j.aim.2016.02.008}{\textit{Adv. Math.}} \textbf{295} (2016), 195--249,
 \href{https://arxiv.org/abs/1212.3264}{arXiv:1212.3264}.

\bibitem{bfk}
Ballard M., Favero D., Katzarkov L., A category of kernels for equivariant
 factorizations and its implications for {H}odge theory, \href{https://doi.org/10.1007/s10240-013-0059-9}{\textit{Publ. Math.
 Inst. Hautes \'Etudes Sci.}} \textbf{120} (2014), 1--111,
 \href{https://arxiv.org/abs/1105.3177}{arXiv:1105.3177}.

\bibitem{bfk2}
Ballard M., Favero D., Katzarkov L., A category of kernels for equivariant
 factorizations, {II}: further implications, \href{https://doi.org/10.1016/j.matpur.2014.02.004}{\textit{J.~Math. Pures Appl.}}
 \textbf{102} (2014), 702--757, \href{https://arxiv.org/abs/1310.2656}{arXiv:1310.2656}.

\bibitem{hl-s}
Halpern-Leistner D., Sam S.V., Combinatorial constructions of derived
 equivalences, \href{https://doi.org/10.1090/jams/940}{\textit{J.~Amer. Math. Soc.}} \textbf{33} (2020), 735--773,
 \href{https://arxiv.org/abs/1601.02030}{arXiv:1601.02030}.

\bibitem{happel}
Happel D., Triangulated categories in the representation theory of
 finite-dimensional algebras, \textit{London Ma\-the\-matical Society Lecture Note
 Series}, Vol.~119, \href{https://doi.org/10.1017/CBO9780511629228}{Cambridge University Press}, Cambridge, 1988.

\bibitem{H2}
Hirano Y., Derived {K}n\"orrer periodicity and {O}rlov's theorem for gauged
 {L}andau--{G}inzburg models, \href{https://doi.org/10.1112/S0010437X16008344}{\textit{Compos. Math.}} \textbf{153} (2017),
 973--1007, \href{https://arxiv.org/abs/1602.04769}{arXiv:1602.04769}.

\bibitem{H1}
Hirano Y., Equivalences of derived factorization categories of gauged
 {L}andau--{G}inzburg models, \href{https://doi.org/10.1016/j.aim.2016.10.023}{\textit{Adv. Math.}} \textbf{306} (2017),
 200--278, \href{https://arxiv.org/abs/1506.00177}{arXiv:1506.00177}.

\bibitem{ks}
Kashiwara M., Schapira P., Sheaves on manifolds, \textit{Grundlehren der
 Mathematischen Wissenschaften}, Vol.~292, \href{https://doi.org/10.1007/978-3-662-02661-8}{Springer-Verlag}, Berlin, 1994.

\bibitem{ks2}
Kashiwara M., Schapira P., Categories and sheaves, \textit{Grundlehren der
 Mathematischen Wissenschaften}, Vol.~332, \href{https://doi.org/10.1007/3-540-27950-4}{Springer-Verlag}, Berlin, 2006.

\bibitem{lip}
Lipman J., Notes on derived functors and {G}rothendieck duality, in Foundations
 of {G}rothendieck duality for diagrams of schemes, \textit{Lecture Notes in
 Math.}, Vol.~1960, \href{https://doi.org/10.1007/978-3-540-85420-3}{Springer}, Berlin, 2009, 1--259.

\bibitem{ls}
Lunts V.A., Schn\"urer O.M., Matrix factorizations and semi-orthogonal
 decompositions for blowing-ups, \href{https://doi.org/10.4171/JNCG/252}{\textit{J.~Noncommut. Geom.}} \textbf{10}
 (2016), 907--979, \href{https://arxiv.org/abs/1212.2670}{arXiv:1212.2670}.

\bibitem{ot}
Okonek C., Teleman A., Graded tilting for gauged {L}andau--{G}inzburg models
 and geometric applications, \href{https://dx.doi.org/10.4310/PAMQ.2021.v17.n1.a5}{\textit{Pure Appl. Math.~Q.}} \textbf{17} (2021),
 185--235, \href{https://arxiv.org/abs/1907.10099}{arXiv:1907.10099}.

\bibitem{olsson}
Olsson M., Algebraic spaces and stacks, \textit{American Mathematical Society
 Colloquium Publications}, Vol.~62, \href{https://doi.org/10.1090/coll/062}{Amer. Math. Soc.}, Providence, RI, 2016.

\bibitem{posi}
Positselski L., Two kinds of derived categories, {K}oszul duality, and
 comodule-contramodule correspondence, \href{https://doi.org/10.1090/S0065-9266-2010-00631-8}{\textit{Mem. Amer. Math. Soc.}}
 \textbf{212} (2011), vi+133~pages, \href{https://arxiv.org/abs/0905.2621}{arXiv:0905.2621}.

\bibitem{rs}
Rennemo J.V., Segal E., Hori-mological projective duality, \href{https://doi.org/10.1215/00127094-2019-0014}{\textit{Duke
 Math.~J.}} \textbf{168} (2019), 2127--2205, \href{https://arxiv.org/abs/1609.04045}{arXiv:1609.04045}.

\bibitem{rsv}
Rennemo J.V., Segal E., Van~den Bergh M., A non-commutative {B}ertini theorem,
 \href{https://doi.org/10.4171/JNCG/334}{\textit{J.~Noncommut. Geom.}} \textbf{13} (2019), 609--616,
 \href{https://arxiv.org/abs/1705.01366}{arXiv:1705.01366}.

\bibitem{svdb}
\v{S}penko \v{S}., Van~den Bergh M., Non-commutative resolutions of quotient
 singularities for reductive groups, \href{https://doi.org/10.1007/s00222-017-0723-7}{\textit{Invent. Math.}} \textbf{210}
 (2017), 3--67, \href{https://arxiv.org/abs/1502.05240}{arXiv:1502.05240}.

\bibitem{stacks}
The Stacks Project Authors, Stacks Project, available at \url{https://stacks.math.columbia.edu}.

\bibitem{tho}
Thomason R.W., Equivariant resolution, linearization, and {H}ilbert's
 fourteenth problem over arbitrary base schemes, \href{https://doi.org/10.1016/0001-8708(87)90016-8}{\textit{Adv. Math.}}
 \textbf{65} (1987), 16--34.

\end{thebibliography}
